% ============================================================================
%  COMP 4905 Honours Project Report
%  Neural Network-Based Cheating Detection for CoMas
%  Sabateesh Sivakumar - Supervisor: Prof. Darryl Hill
%  Carleton University, School of Computer Science
%
%  Compile with:  pdflatex report.tex   (run twice for the table of contents)
% ============================================================================
\documentclass[12pt,oneside]{report}

\usepackage[utf8]{inputenc}
\usepackage[T1]{fontenc}
\usepackage[margin=1in]{geometry}
\usepackage{graphicx}
\usepackage{booktabs}
\usepackage{amsmath}
\usepackage{array}
\usepackage{longtable}
\usepackage{caption}
\usepackage{xcolor}
\usepackage{listings}
\usepackage{enumitem}
\usepackage[hidelinks]{hyperref}

\definecolor{codebg}{RGB}{247,247,249}
\definecolor{codekw}{RGB}{160,32,120}
\definecolor{codecm}{RGB}{100,120,100}
\definecolor{codestr}{RGB}{170,60,40}

\lstdefinestyle{py}{
  language=Python,
  backgroundcolor=\color{codebg},
  basicstyle=\ttfamily\footnotesize,
  keywordstyle=\color{codekw},
  commentstyle=\itshape\color{codecm},
  stringstyle=\color{codestr},
  numbers=left,
  numberstyle=\tiny\color{gray},
  numbersep=8pt,
  frame=single,
  rulecolor=\color{gray!40},
  breaklines=true,
  showstringspaces=false,
  captionpos=b,
  xleftmargin=16pt,
}
\lstset{style=py}

\newcommand{\code}[1]{\texttt{\small #1}}

\setlength{\parskip}{6pt}
\setlength{\parindent}{0pt}
\renewcommand{\arraystretch}{1.2}

\begin{document}

% ===================== TITLE PAGE ==========================================
\begin{titlepage}
\centering
\vspace*{2cm}

{\large Carleton University}\\[0.3cm]
{\large COMP 4905 -- Honours Project -- Summer 2026}\\[4cm]

{\Huge\bfseries Neural Network-Based\\[0.3cm] Cheating Detection for CoMas}\\[1cm]

{\Large A Hybrid Vision and Text Pipeline for\\[0.2cm]
Triaging Online Examination Screenshots}\\[5cm]

{\large Sabateesh Sivakumar}\\[0.2cm]
{\large 101229774}\\[1.5cm]

{\large Supervisor: Professor Darryl Hill}\\[0.2cm]
{\large School of Computer Science}\\[2cm]

{\large September 2026}

\vfill
\end{titlepage}

% ===================== ABSTRACT ============================================
\chapter*{Abstract}
\addcontentsline{toc}{chapter}{Abstract}

CoMas is the proctoring tool used at Carleton University to monitor online
examinations. It periodically captures a screenshot of each student's desktop,
and every one of those images is currently reviewed by hand. A single course
offering can produce many thousands of screenshots, which makes exhaustive
manual review both slow and inconsistent between reviewers.

This project develops a screening layer that ranks screenshots by how likely
they are to show prohibited behaviour, so that a proctor reviews the most
suspicious images first rather than reviewing everything. The system addresses
two concrete detection tasks: identifying whether an AI coding assistant
(GitHub Copilot or Cursor) is active inside Visual Studio Code, and identifying
whether a student has navigated away from their Brightspace quiz.

The central technical finding is that these two tasks require fundamentally
different methods. An assistant chat panel and a Brightspace quiz page are both
identifiable from on-screen text, and a deterministic keyword rule over OCR
output handles them accurately and interpretably. Inline ``ghost text'' --- the
dimmed code an assistant suggests at the cursor --- contains no distinctive
words at all, and is invisible to text extraction. Measured on real screenshots
captured during development, an OCR-only detector recovered just 3 of 14 true
positives. Detecting ghost text therefore requires a visual model, and a
ResNet50 convolutional neural network was fine-tuned for this purpose using
transfer learning from ImageNet.

Because access to real CoMas data could not be obtained within the project
timeline, training data was generated synthetically. Two generators render
labelled screenshots programmatically across multiple editor themes, window
resolutions and layouts. A deliberate design decision, and one that materially
affected the outcome, was to construct \emph{hard negatives}: half of the
negative examples show a non-assistant panel docked in the same screen region
as a chat panel, preventing the network from learning the trivial rule
``panel on the right means cheating.''

On a held-out synthetic set generated with an unseen random seed, the final
model detects assistant chat panels with perfect recall and zero false
positives on browser pages and non-chat side panels. Ghost-text detection is
substantially harder: an ablation isolating this discrimination measured an
AUROC of 0.759 at 512-pixel input resolution, improving to 0.856 at 768 pixels.
The mechanism is quantified in the report --- a ghost-text line survives
downscaling as roughly four pixels of image height, whereas a chat panel
occupies a third of the screen. At the shipped operating point the system
recovers 95.0\% of ghost-text screenshots.

The deliverable is a desktop application that ingests a folder of screenshots,
scores them, and presents a ranked review queue with triage controls and an
evidence export suitable for attaching to an academic integrity case.

% ===================== ACKNOWLEDGEMENTS ====================================
\chapter*{Acknowledgements}
\addcontentsline{toc}{chapter}{Acknowledgements}

I would like to thank my supervisor, Professor Darryl Hill, for his guidance
throughout this project, and in particular for the latitude to narrow the
project's scope once it became clear that a general-purpose ``suspicious
screenshot'' classifier was too vague a target to evaluate rigorously. The
decision to focus on two concrete, well-defined detection tasks is the single
change that most improved the quality of this work.

I would also like to acknowledge the School of Computer Science for the
opportunity to pursue an independent project of this scope, and to note that
the experience of running into --- and having to diagnose --- a long sequence of
practical engineering problems taught me at least as much as the machine
learning content did.

% ===================== CONTENTS ============================================
\tableofcontents
\listoftables

% ===================== 1. INTRODUCTION =====================================
\chapter{Introduction}

\section{Chapter Overview}

This chapter introduces the problem the project addresses and explains why I
thought it was worth solving. Section~\ref{sec:motivation} describes how the
project began and the review burden that motivates it.
Section~\ref{sec:objectives} states the project objectives.
Section~\ref{sec:scope} documents how the scope changed from my original
proposal and why, since that change affects how the rest of the report should
be read. Section~\ref{sec:contributions} summarises the contributions, and
Section~\ref{sec:structure} outlines the remainder of the document.

\section{Project Motivation}
\label{sec:motivation}

This project started as a suggestion from Professor Hill rather than as an idea
of my own, and my first reaction to it was shaped less by the machine learning
involved than by having worked as a teaching assistant for three terms.

I should be precise about my own experience here, because it is the honest
version and it matters for how I approached the design. I have not personally
had to review student screenshots as part of any of my TA duties. What I do
have is three terms of familiarity with the shape of TA work --- the marking
queues, the repetitive passes over hundreds of near-identical submissions, and
a fairly good sense of which parts of the job people dread. When the idea was
first described to me, my immediate thought was not ``this is an interesting
classification problem'' but ``this sounds miserable, and it does not need to
be.'' Sitting down to click through several thousand screenshots, the
overwhelming majority of which show a student quietly working, is exactly the
kind of task that burns people out: high volume, low information density, and
almost entirely negative results. If a tool could remove the images that are
plainly not worth anyone's time and leave only those that genuinely need human
judgement, that struck me as a real improvement in the working lives of
colleagues who do carry this responsibility. That framing --- as something
built for the person doing the reviewing, not simply as a classifier ---
stayed with me for the rest of the project and is visible in a number of design
decisions later in this report.

The underlying numbers explain why the problem exists at all. CoMas captures
desktop screenshots at intervals while a student completes an online
examination, so that a proctor can afterwards inspect what was on screen and
determine whether anything prohibited was visible. In practice the volume
defeats the purpose. A single examination with a few hundred students, each
screenshotted every minute or two across a two-hour sitting, produces tens of
thousands of images. Reviewing all of them carefully is not feasible, so review
becomes shallow, sampled, or both.

Shallow review fails in two directions at once. Genuine misconduct is missed
because the one screenshot that showed it was never looked at closely. At the
same time, enforcement becomes inconsistent: whether a given student is caught
depends on which reviewer happened to examine their images and how much time
that reviewer had that week. The inconsistency is arguably the more serious of
the two problems, because academic integrity processes have to be defensible,
and a process whose outcome depends on reviewer fatigue is difficult to
defend.

The insight the whole project rests on is that the reviewer's task is not
really classification --- it is \emph{prioritisation}. A TA does not need
software that decides whether a student cheated, and I would be uncomfortable
building software that claimed to. What they need is the twenty screenshots
most worth looking at placed in front of them first, so their limited attention
goes where it has the best chance of mattering. The remaining thousands are not
deleted or judged; they are simply moved out of the way until someone has
reason to look.

This reframing has a technical consequence that turns out to matter a great
deal later. A triage system can tolerate a much higher false positive rate than
an automated decision system, because a false positive costs a reviewer a few
seconds and a single keystroke to dismiss, whereas a false negative defeats the
entire purpose of running the system. That asymmetry is not a detail --- it
directly determines the operating point I eventually shipped, and it is
revisited quantitatively in Chapter~\ref{ch:evaluation}.

Finally, the specific behaviours I targeted reflect what cheating in a
programming examination actually looks like now. AI coding assistants such as
GitHub Copilot and Cursor are installed by default in many students'
development environments, and they will readily produce a complete solution to
a typical undergraduate examination question. A student no longer needs to
visit a prohibited website; the prohibited resource is embedded in the editor
they are legitimately entitled to use. That makes this a meaningfully different
problem from the classic proctoring task of spotting an open browser tab, and
it is the problem I found most worth working on.

\section{Project Objectives}
\label{sec:objectives}

I set out to accomplish the following:

\begin{enumerate}[itemsep=2pt]
  \item Acquire and curate a labelled dataset of examination screenshots,
        supplemented by synthetically generated data where real samples prove
        insufficient.
  \item Design and train a binary image classifier distinguishing legitimate
        from suspicious screenshots, using a convolutional neural network with
        transfer learning from a pre-trained backbone.
  \item Evaluate the classifier using precision, recall, F1-score and AUROC,
        with particular attention to minimising false negatives so that genuine
        misconduct is not missed.
  \item Deliver a lightweight inference pipeline that processes a batch of
        screenshots and produces a ranked list of flagged images for manual
        review.
\end{enumerate}

Chapter~\ref{ch:conclusion} assesses the finished system against each of these
objectives, including the one I did not fully meet.

\section{Scope Refinement}
\label{sec:scope}

My original proposal described a general binary classifier separating
``legitimate'' from ``suspicious'' screenshots, with suspicious behaviour
characterised loosely as open browser tabs, messaging applications, or
unauthorised reference material. Early in the implementation phase this proved
to be an unworkable specification, for a reason worth stating plainly:
\emph{``suspicious'' is not a visual category}. A screenshot showing a browser
is suspicious during a closed-book quiz and entirely unremarkable during an
open-book take-home. Without a fixed examination context, no consistent
labelling function exists, and without a consistent labelling function there is
nothing coherent to train against and no defensible way to measure accuracy.

I resisted narrowing the scope for longer than I should have, because the
smaller version felt like a retreat from what I had promised. It was not. The
general version was not more ambitious, only less well defined, and I could not
have produced a single trustworthy number from it.

The project was therefore narrowed to two concrete tasks, each with an
unambiguous ground truth:

\begin{enumerate}[itemsep=2pt]
  \item \textbf{AI assistant detection.} Is GitHub Copilot or Cursor visibly
        active in a Visual Studio Code screenshot, either as a chat panel or as
        an inline ghost-text suggestion?
  \item \textbf{Brightspace tab-leave detection.} Has the student navigated
        away from the Brightspace quiz page they are supposed to be on?
\end{enumerate}

Both questions can be answered by looking at a single screenshot, both admit a
labelling rule that two different people would apply identically, and both
correspond to behaviours that examiners actually care about. The narrowing
traded breadth for the ability to make measurable claims, which for a project
of this length was the correct trade.

One further deviation from the proposal should be recorded. The proposal
specified a web application as the reviewer interface. The delivered interface
is a desktop application built with Tkinter. Functionally these are equivalent
for the stated purpose --- both accept a folder of screenshots and return a
ranked review queue --- but the desktop implementation avoids requiring a
proctor to upload potentially sensitive examination screenshots to a web
service, which on reflection is a meaningful privacy advantage for this
particular application. An earlier FastAPI-based web version was implemented
and then removed once the desktop application superseded it.

\section{Contributions}
\label{sec:contributions}

The main contributions of this work are:

\begin{itemize}[itemsep=2pt]
  \item A hybrid detection pipeline that combines a fine-tuned CNN with an OCR
        keyword detector, together with a quantitative demonstration that
        neither component alone is sufficient.
  \item Two synthetic screenshot generators producing labelled training data
        for the two tasks, including a hard-negative construction that prevents
        a specific and easily-overlooked shortcut solution.
  \item An empirical characterisation of the ghost-text detection problem,
        including an input-resolution ablation that isolates the mechanism
        behind the model's principal weakness.
  \item A reviewer-facing desktop application with ranked review, adjustable
        sensitivity, triage state and evidence export.
  \item A set of measured engineering results of independent interest, most
        notably that na\"ive process-level parallelisation of Tesseract OCR is
        approximately six times \emph{slower} than sequential execution unless
        OpenMP threading is explicitly constrained.
\end{itemize}

\section{Report Structure}
\label{sec:structure}

Chapter~\ref{ch:background} reviews relevant background. Chapter~\ref{ch:design}
describes the system architecture and technology choices.
Chapter~\ref{ch:detection} details the detection methods.
Chapter~\ref{ch:data} covers synthetic data generation.
Chapter~\ref{ch:implementation} discusses implementation and performance
engineering. Chapter~\ref{ch:interface} presents the reviewer interface.
Chapter~\ref{ch:evaluation} reports experimental results.
Chapter~\ref{ch:experience} reflects on the development experience and the
problems encountered. Chapter~\ref{ch:conclusion} concludes and outlines future
work.

% ===================== 2. BACKGROUND =======================================
\chapter{Background and Related Work}
\label{ch:background}

\section{Chapter Overview}

This chapter provides the technical background required to follow the rest of
the report. Section~\ref{sec:proctoring} situates the work within automated
proctoring. Section~\ref{sec:assistants} describes how AI coding assistants
appear on screen, which is the observation the entire detection design rests
on. Sections~\ref{sec:transfer} and~\ref{sec:ocr} cover transfer learning and
optical character recognition respectively.

\section{Automated Proctoring}
\label{sec:proctoring}

Commercial remote proctoring systems typically combine webcam monitoring, screen
capture, and browser lockdown. Lockdown browsers prevent navigation away from
an examination page but cannot govern what happens outside the browser, which
is precisely the gap that matters for a programming examination conducted in a
local development environment. Screen capture is the more general mechanism,
and it is what CoMas uses, but it transfers the burden from prevention to
review.

Automated analysis of proctoring data raises well-known concerns around false
accusation and equity, and these concerns directly shaped the design here. The
system is positioned as a triage aid that ranks screenshots for human review,
not as an automated decision-maker. No screenshot is ever labelled as cheating
by the system alone; the output is an ordering and a numeric score, and a human
makes every determination. The interface reinforces this by requiring the
reviewer to explicitly mark each flagged image as reviewed or dismissed.

\section{AI Coding Assistants On Screen}
\label{sec:assistants}

The design of the VS Code detector follows from a single observation: an AI
coding assistant manifests in two visually and textually distinct ways.

\textbf{The chat panel} is a docked side panel, conventionally on the right,
containing a conversation between the user and the assistant along with a
message composer at the bottom. It is large --- typically 25--35\% of screen
width --- and it contains distinctive interface text. GitHub Copilot's panel
displays strings such as ``GitHub Copilot'' and ``Copilot Chat''; Cursor's
displays ``Plan agent'', ``No matching sessions'' and ``Explore and understand
your code''. Both its size and its text make it comparatively easy to detect.

\textbf{Inline ghost text} is the assistant's autocomplete suggestion, rendered
directly in the editor at the cursor position in a dimmed foreground colour.
The student presses Tab to accept it. This is by far the more common mode of
use during actual coding, and it is much harder to detect. It occupies a single
line of the editor, perhaps 15 pixels tall in a screenshot 1000 pixels high.
Critically, \emph{it contains no distinctive text}: the suggestion is ordinary
source code, indistinguishable in content from code the student wrote. The
words ``Copilot'' or ``Cursor'' appear nowhere. The only signal is visual --- a
run of text rendered in a grey that differs from the surrounding syntax
highlighting.

This asymmetry drives every subsequent design decision. Text-based detection
can find the panel and can never find ghost text. A visual model is not merely
preferable for ghost text; it is the only option.

A complication worth noting is that dimmed grey text is not unique to
assistants. Source code comments are also rendered in a muted colour in most
editor themes. A model that learns ``grey text implies assistant'' would flag
every well-commented file. This informed the construction of the training
corpus described in Chapter~\ref{ch:data}.

\section{Transfer Learning for Image Classification}
\label{sec:transfer}

Training a deep convolutional network from random initialisation requires
datasets far larger than this project could assemble. Transfer learning avoids
this by starting from a network pre-trained on a large general-purpose corpus
--- here ImageNet --- and fine-tuning it on the target task. The early
convolutional layers of such a network encode generic visual primitives (edges,
textures, colour gradients) that transfer readily to new domains, so only the
later layers and classification head need to adapt substantially.

This project uses ResNet50, a 50-layer residual network. Residual connections
allow gradients to propagate through deep stacks without vanishing, which makes
networks of this depth trainable in practice. The final fully-connected
classification layer is replaced with a single-output head producing one logit,
which is passed through a sigmoid to yield a probability of the positive class.

Screenshots differ from ImageNet photographs in ways worth acknowledging: they
are synthetic renderings with hard edges, flat colour regions, and very
fine-grained text detail rather than natural textures. That ImageNet features
transfer usefully at all is not a given, though the empirical results in
Chapter~\ref{ch:evaluation} indicate they do. The resolution ablation in
Section~\ref{sec:ablation} suggests the fine-detail requirement is precisely
where the approach strains.

\section{Optical Character Recognition}
\label{sec:ocr}

OCR converts pixels containing rendered text into machine-readable strings.
This project uses Tesseract, a mature open-source engine, via the
\code{pytesseract} Python binding. Screenshots are an unusually favourable
input for OCR: the text is crisply rendered at high contrast rather than
photographed, so recognition accuracy is high with no preprocessing.

The relevant weakness of keyword-based detection is not accuracy but
\emph{brittleness over time}. The keywords are interface strings chosen by a
vendor, and vendors rename things. A Copilot update that changes ``Tab to
accept'' to different wording silently disables that detection rule with no
error and no warning. This is a genuine maintenance burden and is discussed
further in Section~\ref{sec:limitations}. It is also the main argument for
retaining the CNN as the primary signal: the two methods fail under different
and largely uncorrelated conditions.

% ===================== 3. SYSTEM DESIGN ====================================
\chapter{System Design}
\label{ch:design}

\section{Chapter Overview}

This chapter describes the architecture of the delivered system.
Section~\ref{sec:usecases} states the use cases. Section~\ref{sec:pipeline}
describes the processing pipeline. Section~\ref{sec:tech} lists the
technologies used and justifies the significant choices.
Section~\ref{sec:modules} documents the module structure of the codebase.

\section{Use Cases}
\label{sec:usecases}

The system supports a single primary actor, the proctor or teaching assistant
reviewing examination screenshots.

\textbf{Select a detection mode.} The reviewer chooses among three modes:
\emph{VSCODE Cheating} for programming examinations conducted in an editor,
\emph{Brightspace Cheating} for quizzes taken in the browser, and
\emph{VSCODE + Brightspace Cheating} for examinations involving both. The mode
determines which detectors run and, as described in
Section~\ref{sec:combined}, how their outputs are interpreted.

\textbf{Submit a batch of screenshots.} The reviewer selects a set of image
files through a native file dialog. Analysis runs on a background thread with
progress reporting, so the interface remains responsive.

\textbf{Review a ranked queue.} Flagged screenshots are presented in descending
order of suspicion score. The reviewer pages through them with buttons or arrow
keys, each shown with its score and a plain-language evidence line.

\textbf{Adjust sensitivity.} A threshold slider re-filters results live. The
reviewer can loosen the threshold to see marginal cases or tighten it if the
queue is too noisy, without re-running the analysis.

\textbf{Record triage decisions.} Each screenshot can be marked reviewed or
dismissed, so a reviewer working through a large queue can track their
position.

\textbf{Export evidence.} The reviewer exports a timestamped report folder
containing copies of the flagged images, a CSV of all scores, and a summary
suitable for attaching to an academic integrity case.

\section{Processing Pipeline}
\label{sec:pipeline}

Analysis proceeds in four stages.

\textbf{Stage 1: CNN scoring.} If a trained checkpoint is available and the
mode includes VS Code detection, every screenshot is scored by the CNN. Images
are processed in batches of 16 rather than individually, which is discussed in
Section~\ref{sec:batching}.

\textbf{Stage 2: Selective OCR.} Screenshots the model is already confident
about (score $\geq 0.90$) skip OCR entirely, since the outcome cannot change.
The remainder are OCR'd in parallel across CPU cores. This ordering is
deliberate: OCR is by far the most expensive stage, so it runs only where it
can still affect the result. Section~\ref{sec:skiplogic} explains why this
optimisation is applied only to high-confidence detections and never to
confident \emph{negatives}.

\textbf{Stage 3: Detection.} Each screenshot is scored by the enabled
detectors, reading from the warmed caches. The final score is the maximum
across detectors, so any single strong signal is sufficient to flag an image.

\textbf{Stage 4: Ranking and presentation.} Results are sorted by descending
score and passed to the interface, which filters at the active threshold.

\section{Technologies}
\label{sec:tech}

\begin{table}[h]
\centering
\caption{Technologies used and their roles.}
\begin{tabular}{@{}lll@{}}
\toprule
\textbf{Technology} & \textbf{Version} & \textbf{Role} \\
\midrule
Python & 3.11 & Implementation language \\
PyTorch & 2.x & Model definition, training, inference \\
torchvision & 0.x & ResNet50 backbone and ImageNet weights \\
Tesseract & 4.1 & Optical character recognition \\
pytesseract & 0.3 & Python binding for Tesseract \\
Pillow (PIL) & 10.x & Image loading, rendering, augmentation \\
Tkinter & Tk 8.6 & Desktop user interface \\
scikit-learn & 1.x & Evaluation metrics \\
PyYAML & 6.x & Configuration file parsing \\
pytest & 7.x & Automated testing \\
\bottomrule
\end{tabular}
\end{table}

Two choices merit justification.

\textbf{ResNet50 over a Vision Transformer.} The proposal raised the
possibility of a vision transformer. ResNet50 was chosen because transformers
are substantially more data-hungry, and with a synthetic dataset of 800 images
the risk of a transformer memorising generator artefacts rather than learning
the visual signal was high. ResNet50's stronger inductive bias toward local
spatial features is also a better match for a task where the signal is a
localised region of the screen.

\textbf{Tkinter over a web framework.} Tkinter ships with Python, requires no
server, and keeps examination screenshots on the reviewer's machine. Given that
the input is sensitive student data, avoiding a network upload step is a
genuine advantage rather than merely a convenience. The cost is a less
sophisticated interface toolkit, which required some care to produce an
acceptable result and caused one significant portability problem documented in
Section~\ref{sec:tkbug}.

\section{Module Structure}
\label{sec:modules}

The implementation lives in a single Python package, \code{comas/}, comprising
roughly 15 modules.

\begin{table}[h]
\centering
\caption{Modules of the \code{comas} package.}
\begin{tabular}{@{}p{0.28\textwidth}p{0.64\textwidth}@{}}
\toprule
\textbf{Module} & \textbf{Responsibility} \\
\midrule
\code{copilot.py} & VS Code detector: OCR keywords, CNN scorer, signal fusion \\
\code{brightspace.py} & Brightspace detector: quiz-marker rule, IDE recognition \\
\code{ocr.py} & Tesseract wrapper, content-addressed cache, parallel batching \\
\code{model.py} & ResNet50 and fusion model definitions \\
\code{data.py} & Dataset, transforms, augmentation, session-aware splitting \\
\code{train.py} & Training loop, checkpoint selection \\
\code{evaluate.py} & Metric computation \\
\code{infer.py} & Command-line batch inference \\
\code{config.py} & Typed configuration loaded from YAML \\
\code{gui.py} & Desktop reviewer interface \\
\code{synthetic.py} & Synthetic VS Code screenshot generator \\
\code{synthetic\_brightspace.py} & Synthetic browser and quiz generator \\
\code{variant\_report.py} & Per-variant evaluation tool \\
\code{realpairs.py} & Utilities for organising real screenshot pairs \\
\bottomrule
\end{tabular}
\end{table}

The separation between detection logic and interface is strict: no module in
the detection path imports from \code{gui.py}, and every detector is usable
from the command line independently of the interface. This made the detection
logic testable without instantiating a window, which matters because a
graphical toolkit cannot be exercised in a headless environment.

% ===================== 4. DETECTION METHODS ================================
\chapter{Detection Methods}
\label{ch:detection}

\section{Chapter Overview}

This chapter details how each detector works.
Section~\ref{sec:ocrdetector} describes the OCR keyword detector and
Section~\ref{sec:cnndetector} the CNN. Section~\ref{sec:fusion} explains how
their outputs combine and Section~\ref{sec:skiplogic} the confidence-based
optimisation applied to that combination.
Section~\ref{sec:brightspacedetector} covers the Brightspace detector, whose
design changed substantially during development, and
Section~\ref{sec:combined} the interaction between detectors in combined mode.

\section{VS Code Detector: OCR Keyword Signal}
\label{sec:ocrdetector}

The keyword detector extracts text from a screenshot and searches for strings
characteristic of an assistant interface. Keywords are partitioned into strong
and weak sets. A strong keyword is one that essentially cannot appear except in
an assistant panel; a match scores 0.95. Weak keywords are suggestive but
individually inconclusive; a single match scores 0.45 and two or more score
0.75.

Each keyword also maps to the tool it indicates, so the detector reports
\emph{which} assistant was seen. This distinction proved practically relevant:
the real screenshots captured during development showed Cursor rather than
Copilot, and a detector built only around Copilot's interface strings would
have missed them entirely.

\begin{table}[h]
\centering
\caption{A representative subset of the strong keyword set.}
\begin{tabular}{@{}ll@{}}
\toprule
\textbf{Keyword} & \textbf{Tool indicated} \\
\midrule
\code{github copilot} & Copilot \\
\code{copilot chat} & Copilot \\
\code{tab to accept} & Copilot \\
\code{accept suggestion} & Copilot \\
\code{plan agent} & Cursor \\
\code{no matching sessions} & Cursor \\
\code{explore and understand your code} & Cursor \\
\code{before implementing changes} & Cursor \\
\bottomrule
\end{tabular}
\end{table}

An instructive bug shaped this keyword set. The bare word \code{copilot} was
initially included as a keyword. During testing on real screenshots it produced
false positives on screenshots that showed no assistant at all. The cause was
that this project's own source file is named \code{copilot.py}, and that
filename was visible in the VS Code file explorer sidebar. OCR read it, the
substring matched, and the screenshot was flagged. The fix was to remove the
bare word and require multi-word phrases specific to assistant interfaces.

The episode is worth recording because it illustrates a general hazard of
keyword detection: the detector has no notion of \emph{where} on screen a
string appeared or what interface element it belonged to. It matches text
anywhere in a flattened transcript of the entire screen. A file listing, a
browser bookmark, or a lecture slide mentioning Copilot would all trigger it
equally. Multi-word interface phrases mitigate this but do not eliminate it,
and a regression test now pins the behaviour:

\begin{lstlisting}[caption={Regression test for the filename false positive.}]
def test_filename_copilot_does_not_trigger():
    text = "EXPLORER  __init__.py  copilot.py  data.py  config.yaml"
    strong, weak = find_keywords(text)
    assert strong == []
    assert weak == []
\end{lstlisting}

\section{VS Code Detector: CNN Signal}
\label{sec:cnndetector}

The CNN is a ResNet50 initialised from ImageNet weights, with its final
classification layer replaced by a dropout layer followed by a single-output
linear head. It is trained to distinguish screenshots showing an assistant
--- either as ghost text or a chat panel --- from screenshots that do not.

Training uses binary cross-entropy with logits. Because the dataset contains
200 positives against 600 negatives, the loss is weighted by the ratio of
negative to positive examples, computed automatically as 3.06, so the model is
not rewarded for simply predicting the majority class. Optimisation uses AdamW
with a learning rate of $1\times10^{-4}$, weight decay of $1\times10^{-4}$, and
a cosine annealing schedule across 10 epochs.

The checkpoint retained is not the final epoch but the epoch with the best
validation AUPRC. Area under the precision-recall curve was chosen over
accuracy for selection because it is the metric most sensitive to performance
on the minority positive class, which is the class the system exists to find.
Training runs on Apple Silicon GPU via PyTorch's Metal Performance Shaders
backend, falling back to CUDA or CPU as available.

\begin{lstlisting}[caption={Model construction. \code{use\_ocr} selects between the CNN-only and fusion architectures.}]
class CNNOnlyModel(nn.Module):
    def __init__(self, backbone="resnet50", dropout=0.3, pretrained=True):
        super().__init__()
        self.backbone, feat_dim = _build_cnn_backbone(backbone, pretrained)
        self.head = nn.Sequential(
            nn.Dropout(dropout),
            nn.Linear(feat_dim, 1),
        )

    def forward(self, x, text_emb=None):
        return self.head(self.backbone(x)).squeeze(-1)
\end{lstlisting}

A second architecture, \code{FusionModel}, is implemented and available. It
concatenates the CNN feature vector with a sentence-transformer embedding of
the screenshot's OCR text before the classification head, allowing the network
to learn jointly from appearance and text. It was not used for the final
results: with a synthetic training set whose panel text is deliberately generic
(Section~\ref{sec:leakage}), the text branch carries little information, and the
added dependency and inference cost were not justified. It remains as
infrastructure for future work with real data, where text content would be far
richer.

\section{Signal Fusion}
\label{sec:fusion}

The two signals combine by taking the maximum. If either the CNN or the keyword
detector produces a high score, the screenshot is flagged. This is deliberately
permissive, and follows from the asymmetry discussed in
Section~\ref{sec:motivation}: the cost of a false positive is a few seconds of
reviewer attention, while the cost of a false negative is undetected
misconduct.

Maximum was chosen over a weighted average because the two detectors have
disjoint competencies rather than overlapping opinions. The CNN can see ghost
text and the keyword detector cannot; the keyword detector can identify an
assistant the CNN was never trained on and name it explicitly. Averaging would
let one detector's justified ignorance dilute the other's justified confidence.
Averaging a CNN score of 0.95 with a keyword score of 0.0 yields 0.475 --- a
screenshot with visible ghost text would fall below any reasonable threshold
purely because the text detector, which structurally cannot see ghost text,
did not vote for it.

\section{Confidence-Based OCR Skipping}
\label{sec:skiplogic}

OCR dominates runtime, costing roughly 5.8 seconds per screenshot at native
resolution against approximately 0.3 seconds for CNN inference. Since the CNN
runs first, its output can be used to skip OCR where OCR cannot change the
outcome.

The rule implemented is deliberately \emph{one-sided}. If the CNN scores a
screenshot at 0.90 or above, OCR is skipped: the image will be flagged
regardless, so the additional evidence would not alter the decision. If the CNN
scores \emph{low}, OCR still runs.

The asymmetry is essential and was the subject of careful reasoning during
implementation. It is tempting to skip OCR on confident negatives as well ---
it would roughly double the saving, since most screenshots in a real stream are
negatives. But the CNN was trained to recognise ghost text and chat panels. A
screenshot showing an assistant interface the model has never seen, or one
whose theme falls outside the training distribution, may score low while
containing an unambiguous ``GitHub Copilot'' string in plain text. Allowing a
low CNN score to suppress OCR would convert a certain detection into a silent
miss. The behaviour is pinned by test:

\begin{lstlisting}[caption={The model must never be able to veto a keyword detection.}]
def test_low_model_score_never_suppresses_ocr():
    # the model only learned ghost text, so it must not veto an open chat panel
    cache = _CountingCache("GitHub Copilot")
    det = CopilotDetector(ocr_cache=cache, ml_scorer=lambda p: 0.02)
    ev = det.detect(Path("x.png"))
    assert cache.calls == 1
    assert ev.score == 0.95
\end{lstlisting}

\section{Brightspace Detector}
\label{sec:brightspacedetector}

The Brightspace detector uses no machine learning. It is a deterministic rule
over OCR output, and this is a considered choice rather than an omission: a
quiz page is identifiable by its text, so a rule is accurate, interpretable,
requires no training data, and can be corrected immediately if Brightspace
changes its interface. The CNN is applied where learning is necessary; where a
deterministic rule suffices, a rule is the better engineering choice.

The detector's design changed substantially during development, and the change
is more interesting than the final rule.

\textbf{Initial design (blocklist).} The first implementation searched for
off-task site markers --- \code{chatgpt}, \code{stackoverflow},
\code{google.com/search}, \code{chegg}, \code{quizlet} and similar --- and
flagged screenshots containing them. Quiz markers such as ``Submit Quiz'',
``(1 point)'' and ``elapsed'' were used to identify legitimate quiz pages, with
a \code{split\_view} verdict when both appeared.

This had two failure modes. First, it could only detect sites on the list; a
student consulting a site the author had not thought of was invisible. Second,
and more seriously, it produced false positives on legitimate quiz pages,
because quiz questions themselves mention websites. A question asking which
source is acceptable to cite might contain the word ``Wikipedia''; a question
about research methods might mention Google. The student is doing exactly what
they should be, and the detector flags them.

\textbf{Final design (allowlist).} The rule was inverted. Rather than
enumerating prohibited states, the detector recognises the single legitimate
one:

\[
\text{verdict} =
\begin{cases}
\texttt{on\_quiz}\ (0.05) & \text{if quiz markers are present} \\
\texttt{left\_quiz}\ (0.95) & \text{otherwise}
\end{cases}
\]

During a Brightspace examination the quiz should be on screen. Any capture
without it represents a departure, whether to ChatGPT, a messaging application,
a photo library, or anything else. This catches sites never enumerated, and it
cannot be fooled by quiz text that happens to mention a website, because quiz
markers are decisive.

The off-task keyword list was retained, but demoted from determining the
verdict to annotating it. When a tab-leave is detected and a recognised site
appears in the text, the evidence line reads \code{left\_quiz (stackoverflow)};
otherwise it reads \code{left\_quiz (unrecognized page)}. The reviewer gets the
useful detail without the detector depending on it.

The trade-off is that a student displaying the quiz and a prohibited site
side-by-side in one screenshot is classified as \code{on\_quiz}. This is the
price of not flagging quizzes that mention websites, and it was judged the
correct trade: the split-screen case is still caught by any capture in which
the quiz is scrolled away or covered, whereas the false-positive case would
affect honest students on every question that references a source.

\section{Combined Mode}
\label{sec:combined}

Running both detectors together introduced an interaction bug worth
documenting, since it illustrates how a locally correct rule can be globally
wrong.

Under the allowlist rule, any screenshot without quiz markers is a tab-leave
scoring 0.95. In a combined examination --- a programming test with a
Brightspace component --- a VS Code screenshot contains no quiz markers, so
every legitimate coding screenshot was flagged as having left the quiz. The
detector was behaving exactly as specified, but the specification implicitly
assumed a browser-only examination.

The fix adds context-awareness. In combined mode, a screenshot recognised as a
VS Code window contributes nothing from the Brightspace detector; the VS Code
detector judges it instead. Recognition uses a small set of IDE interface
strings, requiring at least two matches so that a quiz question mentioning a
Python file does not accidentally qualify:

\begin{lstlisting}[caption={IDE recognition used to resolve the cross-detector conflict.}]
IDE_KEYWORDS = [
    "explorer", "debug console", "problems", "output", "terminal",
    "spaces: 4", "utf-8", ".py", "outline", "timeline",
]

def looks_like_ide(text: str) -> bool:
    if not text:
        return False
    low = text.lower()
    return sum(1 for k in IDE_KEYWORDS if k in low) >= 2
\end{lstlisting}

The exemption applies only in combined mode. In Brightspace-only mode an IDE
screenshot receives no special treatment, because in a browser-based
examination an open code editor is itself worth a reviewer's attention.
Verified against the real VS Code screenshots collected during development, all
five were correctly recognised as IDE windows and contributed zero Brightspace
score in combined mode.

% ===================== 5. DATA =============================================
\chapter{Training Data Generation}
\label{ch:data}

\section{Chapter Overview}

This chapter describes how training data was produced.
Section~\ref{sec:dataproblem} explains why synthetic data was necessary.
Section~\ref{sec:generator} describes the VS Code generator and
Section~\ref{sec:hardneg} the hard-negative design, which is the most
consequential decision in the chapter. Section~\ref{sec:leakage} covers
prevention of text leakage and Section~\ref{sec:augment} the domain-gap
augmentations.

\section{The Data Access Problem}
\label{sec:dataproblem}

The proposal identified obtaining labelled data as the project's principal risk,
and that assessment proved correct. Real CoMas screenshots contain student work,
personal information, and in some cases visible identity documents. Access
requires institutional data-sharing arrangements and research ethics approval,
processes with lead times measured in months and largely outside a student's
control. Within the project timeline this data did not become available.

The proposal anticipated this and specified a fallback: generating synthetic
data by staging deliberately non-compliant desktop environments. That fallback
became the primary strategy. Three sources of data were ultimately used:

\begin{enumerate}[itemsep=2pt]
  \item \textbf{Synthetic screenshots}, rendered programmatically. These form
        the training and validation sets. Roughly 800 VS Code images and 400
        browser images.
  \item \textbf{Real staged screenshots}, captured manually on the author's own
        machine with Cursor active and inactive. Approximately 40 images, used
        exclusively for qualitative validation --- never for training.
  \item \textbf{Held-out synthetic screenshots}, generated after training with
        an unseen random seed, used for the quantitative evaluation in
        Chapter~\ref{ch:evaluation}.
\end{enumerate}

The limitation this imposes on the results is significant and is stated
explicitly in Section~\ref{sec:validity} rather than left for a reader to
infer.

\section{Synthetic VS Code Screenshot Generator}
\label{sec:generator}

The generator, \code{synthetic.py}, renders VS Code screenshots directly with
Pillow's drawing primitives rather than automating a real editor. Programmatic
rendering gives exact ground truth --- the generator knows precisely which
pixels are ghost text --- and allows systematic variation across conditions
that would be tedious to stage manually.

Each generated screenshot varies along several dimensions: four editor themes
(Dark+, One Dark, Monokai, Light+), four window resolutions from
$1280\times800$ to $1680\times1050$, font size, the presence and width of the
file explorer sidebar, the number and arrangement of open editor tabs, cursor
position, and scroll offset.

The code displayed in the editor is drawn from the project's own Python source
files and tokenised for syntax highlighting with a simplified lexer covering
keywords, strings, numbers, comments and function names. Using real source code
rather than generated text keeps the visual statistics realistic --- genuine
indentation patterns, line-length distributions and identifier shapes.

Comments are deliberately retained in the corpus. Since comments render in a
muted colour similar to ghost text, removing them would allow the network to
learn the shortcut ``any grey text implies assistant''. Their presence forces
it to distinguish ghost text from comments on finer cues: position relative to
the cursor, the absence of a comment marker, and continuation of a code
statement rather than prose.

Positive examples come in three variants: \emph{ghost} (an inline suggestion
only), \emph{panel} (a chat panel only), and \emph{both}. The chat panel is
rendered as a docked right-hand region with a header, alternating user and
assistant message bubbles, occasional embedded code blocks, and a composer box
at the bottom.

\section{Hard Negatives}
\label{sec:hardneg}

The most consequential design decision in the dataset concerns what the
negative examples look like.

The na\"ive approach is to pair chat-panel positives with clean-editor
negatives. The problem is that this makes the classes separable by a feature
that has nothing to do with assistants: \emph{the presence of a docked panel on
the right}. VS Code has many such panels --- Outline, Extensions, Source
Control, Search, split editors. A model trained against clean editors alone
learns ``panel on the right implies cheating'' because that rule achieves
perfect training accuracy, and would then flag every student with the
Extensions view open.

To prevent this, half of the negative examples render a non-assistant panel in
the same screen region, at the same range of widths, with the same header
styling. The model cannot use position or the mere existence of a panel; it
must learn the chat panel's internal structure --- stacked message bubbles
above a composer box --- to separate the classes.

A second category of negative was added later, after the first training run.
The initial dataset's negatives were all VS Code editors, meaning the model had
never seen a screenshot that was not an IDE. A real CoMas capture stream is
mostly browsers, document editors and desktops. A model whose entire experience
is editors has undefined behaviour on a browser window, and could plausibly
flag every Brightspace quiz page as an assistant. The 400 synthetic browser
screenshots from the Brightspace generator --- all assistant-free by
construction --- were therefore folded in as additional negatives.

\begin{table}[h]
\centering
\caption{Composition of the VS Code training set (\code{data\_vscode}).}
\begin{tabular}{@{}llrl@{}}
\toprule
\textbf{Class} & \textbf{Variant} & \textbf{Count} & \textbf{Content} \\
\midrule
Positive & ghost & 80 & Inline suggestion at cursor \\
Positive & panel & 80 & Assistant chat panel docked right \\
Positive & both & 40 & Ghost text and panel together \\
\midrule
Negative & clean & 100 & Editor, no panel, no ghost text \\
Negative & hardneg & 100 & Non-assistant panel docked right \\
Negative & browser & 400 & Quiz and non-quiz browser pages \\
\midrule
\multicolumn{2}{@{}l}{\textbf{Total}} & \textbf{800} & \\
\bottomrule
\end{tabular}
\end{table}

The variant of each image is encoded in its filename, which enables the
per-variant analysis in Section~\ref{sec:variant}. This was added after an
earlier training run produced a single aggregate accuracy figure from which it
was impossible to determine whether the model detected panels, ghost text, or
only one of the two. Aggregate metrics conceal exactly the failure that matters
most here.

\section{Preventing Text Leakage}
\label{sec:leakage}

A subtle hazard in generating synthetic data for a vision model is that the
rendered text may make the task trivially solvable by reading rather than
seeing. If the synthetic chat panels displayed ``GitHub Copilot'' in their
headers, a fusion model's text branch could separate the classes perfectly
without the CNN learning anything visual, and reported accuracy would say
nothing about ghost-text detection.

Two measures prevent this. Source files whose names or contents contain
\code{copilot}, \code{cursor} or \code{ghost} are excluded from the code corpus,
and matching lines are filtered out. Chat panel text is deliberately generic:
prompts are ordinary programming questions, replies are ordinary explanations,
and the composer placeholder reads ``Ask anything'' rather than any
vendor-specific string.

This was verified empirically. Running the OCR keyword detector across a sample
of the generated dataset flagged 0 of 15 positives and 0 of 15 negatives,
confirming that no keyword signal is available and the CNN must rely on visual
structure. A consequence worth stating is that these synthetic images are
useless for evaluating the OCR detector --- that evaluation requires real
screenshots, and is reported separately in Section~\ref{sec:ocrbaseline}.

\section{Domain-Gap Augmentation}
\label{sec:augment}

Synthetic renders are pixel-perfect. Real screenshots have been captured from a
high-DPI display, possibly rescaled, and possibly compressed. A model trained
exclusively on crisp renders may key on that crispness.

Three train-time augmentations narrow this gap: random downscale-then-upscale
(simulating resolution loss), random JPEG compression at quality 40--90
(introducing block artefacts), and slight Gaussian blur. Each applies
probabilistically, so the model sees both degraded and clean versions across
epochs.

These are implemented as plain classes rather than lambdas for a specific
practical reason. PyTorch's \code{DataLoader} with multiple worker processes
must serialise transforms to those workers, and on macOS workers are spawned
rather than forked, so transforms must be picklable. A lambda is not picklable
and would crash training with a multi-worker configuration. A test asserts
picklability so this cannot silently regress.

An honest note on these augmentations: they narrow the domain gap in one
direction but also actively suppress the ghost-text signal, since a one-pixel
line of dimmed text is exactly what blur and compression destroy first. The
resolution ablation in Section~\ref{sec:ablation} should be read with this
interaction in mind.

% ===================== 6. IMPLEMENTATION ===================================
\chapter{Implementation and Performance}
\label{ch:implementation}

\section{Chapter Overview}

This chapter covers implementation decisions that materially affected
correctness or performance. Section~\ref{sec:cache} describes the OCR cache.
Section~\ref{sec:parallel} reports the parallelisation work, including a result
that inverted the expected outcome. Section~\ref{sec:roi} covers region-of-
interest OCR, Section~\ref{sec:batching} batched inference, and
Section~\ref{sec:splitting} session-aware data splitting.
Section~\ref{sec:testing} describes the testing strategy.

\section{OCR Caching}
\label{sec:cache}

OCR is expensive and its result depends only on image content, making it an
ideal caching target. The cache is keyed by SHA-1 hash of the file's bytes
rather than by path, so identical images cached under one filename are found
under another, and modified files miss correctly.

Two robustness properties were added after considering realistic failure. The
cache flushes to disk every 25 entries rather than only at completion: a crash
or a closed window at image 190 of 200 would otherwise discard all completed
work. Writes go to a temporary file followed by an atomic rename, so an
interruption mid-write cannot corrupt the cache and lose every previously
cached result.

Region crops are cached under distinct keys from full-frame extractions, so
that a partial-image transcript can never satisfy a later full-image lookup.

\section{Parallel OCR and an Unexpected Result}
\label{sec:parallel}

Tesseract is single-threaded per image and OCR dominates runtime, so
distributing images across a process pool appeared to be a straightforward
optimisation. Measured on a 4-core machine, the first implementation processed
20 images in 29.6 seconds against 4.6 seconds sequentially --- approximately
\textbf{six times slower}.

The cause is OpenMP. Tesseract 4 is internally multi-threaded, and each worker
process spawned its own thread pool. With several workers each spawning threads
on a 4-core machine, the resulting oversubscription caused contention that
overwhelmed any benefit from parallelism. Constraining each worker to a single
OpenMP thread resolved it. The constraint is applied inside a pool initialiser
so that it is automatic rather than dependent on the user's environment:

\begin{lstlisting}[caption={Without this initialiser, parallelisation is a large performance regression.}]
def _worker_init():
    # Tesseract multithreads through OpenMP. Left alone, every worker spawns
    # its own thread pool and they fight over the cores - measurably slower
    # than running sequentially.
    os.environ["OMP_THREAD_LIMIT"] = "1"
\end{lstlisting}

\begin{table}[h]
\centering
\caption{OCR parallelisation, 20 synthetic images, 4 cores, 3 workers.}
\begin{tabular}{@{}lrr@{}}
\toprule
\textbf{Configuration} & \textbf{Time (s)} & \textbf{Speedup} \\
\midrule
Sequential & 4.6 & 1.00$\times$ \\
Parallel, unconstrained OpenMP & 29.6 & 0.16$\times$ \\
Parallel, one thread per worker & 2.6 & 1.79$\times$ \\
Parallel + region-of-interest & 2.2 & 2.13$\times$ \\
\bottomrule
\end{tabular}
\end{table}

On real screenshots at $3600\times2338$, six images took 17.6 seconds
sequentially and 12.4 seconds in parallel, a more modest $1.42\times$. The
smaller gain reflects memory bandwidth limits when several workers each hold a
large image. The general lesson --- that a plausible optimisation can be a
substantial regression, and that measurement rather than reasoning settles the
question --- recurred often enough in this project to be worth stating once
plainly.

A related measurement informed a decision not taken. Downscaling screenshots
before OCR reduced per-image time from 5.84 to 3.09 seconds, an attractive
$1.89\times$. But on real screenshots the keyword verdict changed on 2 of 6
images, including one true positive dropping from 0.95 to 0.00, because small
interface text ceased to be legible. The optimisation was rejected. It would
have been easy to adopt on the strength of the timing number alone without
checking whether the output was still correct.

\section{Region-of-Interest OCR}
\label{sec:roi}

For Brightspace detection, the decisive evidence --- tab titles and the URL bar
--- sits in the top portion of the screen. OCR'ing only the top 25\% is
correspondingly faster.

Measured directly, cropping alone was not safe: it disagreed with full-frame
analysis on 2 of 12 images, in each case downgrading a genuine tab-leave
because that particular page had nothing recognisable in its top strip. The
failure mode was specific --- the crop found nothing \emph{at all} --- which
made it escapable. The implemented approach is two-stage: OCR the top strip
first, and if it yields no quiz markers, re-read the full frame before
concluding anything. Verdicts then matched full-frame analysis on 16 of 16
test images while retaining most of the speed benefit.

\section{Batched CNN Inference}
\label{sec:batching}

The interface originally scored screenshots one at a time in a Python loop,
performing a full forward pass on a batch of one. Since the fixed overhead of a
forward pass is substantial relative to the marginal cost of additional images
in a batch, images are now stacked into batches of 16.

The scorer remains callable per-image for compatibility with the command-line
tools, but the interface uses \code{score\_batch}. An unreadable or corrupt file
scores 0.0 and logs a warning rather than raising, so one bad file in a folder
of 500 cannot abort an entire review session.

\section{Session-Aware Splitting}
\label{sec:splitting}

Screenshots from the same examination session are highly correlated: same
theme, same window size, same file open, often consecutive frames differing by
a few lines of scroll. Splitting such data randomly leaks information between
training and test sets, because a test image may be nearly identical to a
training image, and reported accuracy becomes optimistic.

Splitting is therefore performed at the session level. Each filename begins
with a session identifier, and all images from a session are assigned wholly to
one split. The generator distributes images across 20 synthetic sessions for
this purpose.

This principle initially had a gap in the evaluation tooling. The per-variant
report tool originally scored every image in the dataset, including the 560 the
model had trained on, inflating every figure it produced. The tool now
reproduces the exact split used during training --- same seed, same fractions
--- and defaults to reporting only held-out test sessions.

\section{Testing}
\label{sec:testing}

The project has 42 automated tests covering keyword matching, scoring rules,
cache behaviour, augmentation properties, the synthetic generator's label
balance, and interface logic.

Two testing decisions are worth noting. Pure logic was deliberately separated
from interface code as static methods, so that scoring and evidence-formatting
rules can be tested without instantiating a window --- graphical toolkits
cannot be exercised in a headless environment. And every behavioural bug found
during development had a regression test added alongside its fix, including the
\code{copilot.py} filename false positive, the one-sided OCR skip rule, and the
combined-mode IDE exemption. These tests encode decisions whose rationale is
not obvious from the code alone, and which a future maintainer might otherwise
reasonably reverse.

% ===================== 7. INTERFACE ========================================
\chapter{The Reviewer Interface}
\label{ch:interface}

\section{Chapter Overview}

This chapter describes the delivered application. Section~\ref{sec:uigoals}
states the design goals, Section~\ref{sec:uiwalk} walks through the workflow,
and Section~\ref{sec:uifeatures} details the review features.

\section{Design Goals}
\label{sec:uigoals}

The interface was designed around three principles.

\textbf{The reviewer decides.} The system never labels a screenshot as
cheating. It produces a score and an ordering, and every determination is made
by a human. This is reflected in the vocabulary --- screenshots are ``flagged''
for review, not identified as violations --- and in requiring explicit triage
decisions.

\textbf{Sensitivity is the reviewer's choice.} Different courses have different
tolerances, and the appropriate trade-off between missed detections and
reviewer time is a policy question rather than a technical one. The threshold
is exposed as a live control rather than fixed in code.

\textbf{Institutional visual identity.} The application follows Carleton
University's brand: the official wordmark, Carleton red (\code{\#E91C24}), and
a flat, clean layout. For a tool intended for administrative use, looking like
institutional software rather than a student project is not a trivial concern.

\section{Workflow}
\label{sec:uiwalk}

\textbf{Mode selection.} On launch the reviewer chooses one of three detection
modes. The upload control remains disabled until a mode is chosen, making the
required order unambiguous.

\textbf{Upload and analysis.} A native file dialog accepts multiple images.
Analysis runs on a background thread so the window stays responsive, and the
status line reports each stage --- model scoring, text extraction with the
number of cores in use, and final analysis --- with running counts.

\textbf{Review.} Flagged screenshots appear in descending score order. Each is
displayed on a dark stage that expands to fill available window space and
re-fits when the window is resized. Images are resampled with the Lanczos
filter, which preserves the legibility of small interface text --- important
when the evidence under review may be a single line of dimmed code.

\section{Review Features}
\label{sec:uifeatures}

\textbf{Threshold slider.} Adjusting the threshold re-filters results
immediately without re-running detection. The implementation stores raw
per-detector scores rather than pre-computed verdicts, and derives displayed
evidence at render time, so the evidence line stays consistent with the active
threshold.

\textbf{Triage state.} Each screenshot can be marked reviewed or dismissed,
displayed as a badge and carried into the export. Keyboard shortcuts
(\code{r}, \code{d}, \code{c}, arrow keys) allow a reviewer to work through a
long queue without the mouse.

\textbf{Evidence export.} Export produces a timestamped folder containing
copies of flagged images excluding dismissed ones, a CSV recording every
screenshot with per-detector scores, flag state, triage decision and evidence
text, and a summary file recording the mode, threshold and counts. Recording
the threshold matters for auditability: a reviewer's decisions are only
interpretable alongside the sensitivity at which they were made.

% ===================== 8. EVALUATION =======================================
\chapter{Evaluation and Results}
\label{ch:evaluation}

\section{Chapter Overview}

This chapter reports experimental results. Section~\ref{sec:setup} describes
the setup and Section~\ref{sec:metrics} the metrics.
Section~\ref{sec:ocrbaseline} establishes the OCR baseline on real screenshots
--- the measurement that motivates the CNN. Section~\ref{sec:cnnresults}
reports overall CNN performance and Section~\ref{sec:variant} the per-variant
breakdown that reveals where it actually succeeds.
Section~\ref{sec:ablation} presents the resolution ablation and
Section~\ref{sec:operating} the operating point selection.
Section~\ref{sec:validity} states threats to validity.

\section{Experimental Setup}
\label{sec:setup}

All experiments ran on an Apple Silicon MacBook Pro using PyTorch's Metal
Performance Shaders backend. Training used 10 epochs at batch size 8 for
768-pixel input, approximately five minutes per epoch.

Three datasets are used. \textbf{Training} (\code{data\_vscode}): 800 images,
split 560 train / 150 validation / 90 test by session.
\textbf{Held-out synthetic} (\code{data\_vscode\_holdout}): 255 images generated
after training with random seed 777, so no image and no session was seen during
training --- 100 positives (40 ghost, 40 panel, 20 both) and 155 negatives
(50 clean, 50 hard-negative, 55 browser). \textbf{Real screenshots}:
approximately 40 images captured manually with Cursor active or absent, used
for the OCR baseline and qualitative validation.

\section{Metrics}
\label{sec:metrics}

Recall is emphasised throughout, per the objective of minimising false
negatives. AUROC and AUPRC are reported as threshold-free summaries, with AUPRC
preferred for checkpoint selection because it is more sensitive to
minority-class performance.

One additional metric was defined for this project. \textbf{Ghost-vs-clean
AUROC} is computed using only ghost-text positives and clean-editor negatives,
excluding panels and browsers. Overall AUROC is dominated by easy comparisons
--- a chat panel against a browser page is trivially separable --- and conceals
performance on the hardest discrimination, which is a dimmed line of code
against a plain editor. This metric isolates it.

\section{OCR Baseline on Real Screenshots}
\label{sec:ocrbaseline}

The keyword detector was evaluated on the real staged screenshots. Of 14 real
positives, it flagged 3 --- a recall of 0.214. On a further 5 real positives
held out separately, it flagged 0.

Inspection of the failures confirmed the predicted mechanism. The three
detections all had an assistant chat panel or tooltip open, supplying matching
interface text. The eleven misses showed inline ghost text with no panel: a
dimmed line of suggested code beside a small gutter icon, with no assistant
text anywhere on screen. The OCR transcript of these images contains only the
student's source code, which is indistinguishable from any other code.

This measurement is the empirical justification for the CNN. It also indicates
that a keyword-only system --- which would be far simpler to build and
maintain --- would miss roughly four out of five real cases.

\section{CNN Results}
\label{sec:cnnresults}

Three training runs were conducted.

\begin{table}[h]
\centering
\caption{Training runs. Run 1 predates the addition of browser negatives.}
\begin{tabular}{@{}llrrrr@{}}
\toprule
\textbf{Run} & \textbf{Res.} & \textbf{Images} & \textbf{Test AUROC} & \textbf{Test AUPRC} & \textbf{Acc@0.5} \\
\midrule
1 & 512 & 400 & 0.941 & 0.956 & 0.900 \\
2 & 512 & 800 & 0.996 & 0.958 & 0.967 \\
3 & 768 & 800 & 0.994 & 0.944 & 0.978 \\
\bottomrule
\end{tabular}
\end{table}

These aggregate figures are reported for completeness but should be treated
with caution, for two reasons. The test split contains only 90 images, of which
few are ghost-text positives, so aggregate metrics are dominated by easy
negatives. More importantly, aggregate accuracy conceals the distinction that
matters --- whether the model detects ghost text or only panels. That question
requires the per-variant analysis below.

\section{Per-Variant Analysis}
\label{sec:variant}

Evaluating the 768-pixel model on the held-out synthetic set, broken down by
variant, gives a much clearer picture.

\begin{table}[h]
\centering
\caption{Per-variant performance, 768-pixel model, held-out set, threshold 0.50.}
\begin{tabular}{@{}llrr@{}}
\toprule
\textbf{Class} & \textbf{Variant} & \textbf{Count} & \textbf{Recall / FPR} \\
\midrule
Positive & panel & 40/40 & recall 1.000 \\
Positive & both & 20/20 & recall 1.000 \\
Positive & ghost & 28/40 & recall 0.700 \\
\midrule
Negative & hardneg & 0/50 & FPR 0.000 \\
Negative & browser & 0/55 & FPR 0.000 \\
Negative & clean & 5/50 & FPR 0.100 \\
\bottomrule
\end{tabular}
\end{table}

The system has effectively two distinct competencies rather than one.

\textbf{Panel detection is solved} on this data: perfect recall on both
panel-only and combined positives, with zero false positives on hard negatives
and browser pages. The hard-negative construction succeeded --- the model
distinguishes an assistant chat panel from an Outline or Extensions panel in
the same screen position without error, confirming it learned internal
structure rather than position.

\textbf{Ghost-text detection is materially weaker}, at 0.700 recall. Every
false positive in the entire evaluation occurred on clean editors --- the class
that differs from ghost-text positives by a single dimmed line. The
discrimination the system finds hard is precisely the one predicted to be hard
in Section~\ref{sec:assistants}.

\section{Input Resolution Ablation}
\label{sec:ablation}

To test the hypothesis that ghost-text difficulty is driven by input
resolution, an ablation was run changing only \code{img\_size} from 512 to 768,
with identical data, architecture, schedule and seed.

\begin{table}[h]
\centering
\caption{Resolution ablation on the held-out set. Only input resolution differs.}
\begin{tabular}{@{}lrr@{}}
\toprule
\textbf{Metric (threshold 0.50)} & \textbf{512 px} & \textbf{768 px} \\
\midrule
Ghost-vs-clean AUROC & 0.759 & \textbf{0.856} \\
Ghost recall & 0.825 & 0.700 \\
Clean-editor FPR & 0.540 & \textbf{0.100} \\
Panel recall & 1.000 & 1.000 \\
Hard-negative FPR & 0.000 & 0.000 \\
\bottomrule
\end{tabular}
\end{table}

Ghost-vs-clean AUROC improves from 0.759 to 0.856. The recall and FPR columns
must be read together to avoid a misleading conclusion: the 512-pixel model
appears to have higher ghost recall, but achieves it while flagging 54\% of
clean editors. Its scores are simply shifted, so a fixed threshold of 0.50 sits
at a different point on its curve. The threshold-free AUROC shows the
768-pixel model is better at every operating point --- it attains comparable
recall at a fifth of the false positive rate.

The mechanism is arithmetic. A ghost-text line rendered at 12--15 pixels in a
$1680\times1050$ screenshot survives downscaling to 512 pixels as roughly four
pixels of image height, before the augmentations of Section~\ref{sec:augment}
further blur and compress it. A chat panel occupying a third of the screen
survives any downscaling comfortably. Higher input resolution preserves more of
the only signal that distinguishes ghost text from a clean editor.

That 0.856 remains well short of the 1.000 achieved on panels indicates
resolution is a contributing factor rather than the whole explanation, and
motivates the tiling approach proposed in Section~\ref{sec:future}.

\section{Operating Point Selection}
\label{sec:operating}

Because the model ranks ghost text well but separates it imperfectly, the
threshold choice materially changes system behaviour.

\begin{table}[h]
\centering
\caption{Threshold sweep, 768-pixel model, held-out set. Panel recall is 1.000 and hard-negative and browser FPR are 0.000 at all three thresholds.}
\begin{tabular}{@{}lrr@{}}
\toprule
\textbf{Threshold} & \textbf{Ghost recall} & \textbf{Clean-editor FPR} \\
\midrule
0.50 & 0.700 & 0.100 \\
0.35 & 0.825 & 0.380 \\
\textbf{0.20 (shipped)} & \textbf{0.950} & 0.820 \\
\bottomrule
\end{tabular}
\end{table}

The shipped default is 0.20, chosen deliberately. The project objectives
specify particular attention to minimising false negatives, and the review
workflow makes the trade favourable: results are presented in descending score
order and a reviewer dismisses a benign flag with a single keystroke, whereas a
missed detection is invisible and unrecoverable. A reviewer who finds the queue
too noisy can raise the threshold; they cannot recover a screenshot that was
never surfaced.

What an 0.82 clean-editor false positive rate means in practice should be
stated plainly rather than glossed. At this threshold the system is no longer
\emph{filtering} plain editor screenshots --- it surfaces most of them. Its
contribution for that class is \emph{ranking}: genuine ghost text scores higher
and sorts to the top of the queue. The structural negatives behave quite
differently, with browsers, quiz pages and non-chat panels correctly excluded
at 0\% false positives across all thresholds.

An accurate one-sentence characterisation of the delivered system is therefore:
it eliminates browser pages, quiz pages and non-assistant panels entirely,
detects assistant chat panels with perfect accuracy, and ranks plain editor
screenshots by ghost-text likelihood for human review.

\section{Threats to Validity}
\label{sec:validity}

\textbf{Synthetic evaluation.} The quantitative results are measured on
synthetic data. Although the held-out set uses an unseen seed and unseen
sessions, it shares a generator with the training data --- the same drawing
code, the same four themes, the same fonts. Real screenshots differ in ways the
generator does not model: unusual themes, extensions altering the interface,
partial renders, non-standard window sizes. The reported figures should be read
as an upper bound on real-world performance.

\textbf{Small real sample.} The OCR baseline rests on 19 real positives. The
direction of that result is unambiguous --- a detector that structurally cannot
see ghost text will not see it --- but the precise recall figure of 0.214 has
wide error bars.

\textbf{No real-data CNN evaluation.} The CNN has not been evaluated
quantitatively on real screenshots, because no labelled real set of sufficient
size exists. This is the most significant limitation of the evaluation and
follows directly from the data access problem in
Section~\ref{sec:dataproblem}.

\textbf{Single architecture.} Only ResNet50 was evaluated. No comparison
against other backbones or a vision transformer was performed.

% ===================== 9. EXPERIENCE =======================================
\chapter{Development Experience}
\label{ch:experience}

\section{Chapter Overview}

This chapter reflects on the process of building the system rather than its
final state. Section~\ref{sec:whatworked} describes approaches that proved
valuable. Section~\ref{sec:problems} documents the significant problems
encountered and how they were diagnosed, and Section~\ref{sec:differently}
states what I would do differently.

\section{What Worked}
\label{sec:whatworked}

\textbf{Narrowing the scope.} The most valuable decision was abandoning the
general ``suspicious screenshot'' framing for two specific tasks. It converted
a project I could not evaluate into one I could measure, and every subsequent
result in this report depends on it. I initially resisted the change because
the narrower version felt less ambitious; in retrospect the general version was
not more ambitious, merely less well defined.

\textbf{Measuring instead of assuming.} Several plausible optimisations turned
out to be wrong, and only measurement revealed it. Parallelising OCR made it
six times slower. Downscaling before OCR was nearly twice as fast and silently
broke detections. Cropping to a region of interest was faster and lost true
positives. In each case my reasoning was sound and the outcome was the
opposite. I now treat ``this should be faster'' as a hypothesis rather than a
conclusion.

\textbf{Building diagnostics rather than reading aggregate numbers.} The
per-variant report was the single most useful tool I built. Before it existed I
had a model with 0.94 AUROC and no idea whether it detected ghost text at all.
Afterwards I could see that panel detection was perfect and ghost detection was
at 0.70, which reframed the entire remaining effort. Aggregate metrics told me
the model was good; the breakdown told me what it was good at.

\textbf{Regression tests for reasoning, not just code.} Several tests exist to
preserve decisions whose rationale is invisible in the code --- that a low CNN
score must never suppress OCR, that a filename must not trigger a keyword. Each
encodes a subtle argument that a future reader, including me, could easily
undo.

\section{Problems Encountered}
\label{sec:problems}

\subsection{A blank application window}
\label{sec:tkbug}

The first run of the interface produced a window in which the header and title
simply did not render, showing content bleeding through from behind. The layout
code was correct and the same code behaved correctly elsewhere.

The cause was the Python interpreter. macOS ships an Xcode Python 3.9 bundled
with Tk 8.5.9, a version with known rendering defects on modern macOS; widgets
constructed alongside a \code{PhotoImage} load intermittently fail to paint.
Rebuilding the virtual environment against Homebrew Python 3.11 with Tk 8.6
resolved it completely.

The lesson was about diagnosis rather than Tk. I spent time inspecting my
layout code because that was what I had written and therefore what I assumed
was wrong. The evidence --- content from another window showing through, which
no layout bug can cause --- pointed at the rendering layer from the beginning.
I have since recorded the interpreter requirement in the project README so the
problem cannot recur silently for anyone else.

\subsection{A repository that could not be pushed}

Several weeks in, pushing to GitHub failed: a file exceeded the 100 MB limit.
The file was \code{libtorch\_cpu.dylib}, part of PyTorch, inside a virtual
environment that had been committed before a \code{.gitignore} existed.

Because the offending commits had never been pushed, the fix was to reset to
the last shared commit and recreate the work as a single clean commit
respecting \code{.gitignore}. Adding the directory to \code{.gitignore} alone
would not have helped, since the file remained in history.

The underlying mistake was committing before thinking about what belonged under
version control. I subsequently accumulated four virtual environments totalling
2.4 GB, three of them created while diagnosing the Tk problem and never
removed.

\subsection{A detector that detected itself}

Testing the keyword detector on real screenshots produced false positives on
images containing no assistant. The cause was that my own source file is named
\code{copilot.py}, it was visible in the editor's file explorer, and the bare
keyword \code{copilot} matched it.

This is my favourite bug in the project because the detector was working
perfectly --- it found the string ``copilot'' on screen, which is exactly what
it was told to do. The specification was wrong, not the implementation. It
clarified something about keyword detection generally: the detector has no
concept of \emph{where} text appeared or what interface element it belonged to,
and any such system will match text in file listings, bookmarks and documents
with equal enthusiasm.

\subsection{Correct rules that were globally wrong}

Two bugs shared a structure worth naming. The Brightspace detector flagged
every VS Code screenshot in combined mode, because a screenshot with no quiz
markers is a tab-leave --- correct in a browser-only examination, wrong when
the student is legitimately in an editor. Separately, the off-task keyword list
flagged legitimate quiz pages whose questions mentioned Wikipedia or Google.

In both cases the rule was locally correct and the surrounding assumption was
unstated. This is the class of bug I found hardest to anticipate, because
testing a component in isolation cannot reveal it. Both were found by running
the full application on realistic input rather than by unit testing.

\subsection{Evaluating on training data}

The first version of the per-variant report scored every image in the dataset,
including the 560 the model had trained on. The numbers it produced were
inflated and, worse, plausible --- nothing about the output indicated anything
was wrong.

The fix was to have the tool reproduce the training split exactly and report
only held-out sessions by default. That an evaluation tool can silently produce
meaningless-but-believable numbers is a lesson I expect to carry well beyond
this project.

\subsection{Data access}

The problem with the largest effect on the final result is the one I made least
progress on. The proposal identified real CoMas data as the primary source and
correctly named it as the project's principal risk. I did not begin the access
process early enough, and the timeline for institutional data agreements and
ethics review is not something a student can compress.

The synthetic fallback was anticipated in the proposal and does work, but every
number in Chapter~\ref{ch:evaluation} carries an asterisk as a result. If I
were starting again, the access request would be the first task of week one,
before any code was written, precisely because its lead time is long and
entirely outside my control. Risks that cannot be compressed later must be
started first.

\section{What I Would Do Differently}
\label{sec:differently}

\textbf{Start the data request immediately.} As above --- this is the change
that would most improve the work.

\textbf{Build the per-variant evaluation before the model.} I trained for
several runs while only able to see a single aggregate number. Had the
breakdown existed from the start, I would have identified the ghost-text
weakness far earlier and had more time to address it.

\textbf{Question resolution earlier.} The input size of 512 pixels was carried
forward without examination for most of the project. It is the single
hyperparameter most directly implicated in the model's main weakness, and one
afternoon's ablation produced the clearest result in the report. I should have
asked what the signal physically looks like after preprocessing much sooner.

\textbf{Set up version control properly on day one.} A \code{.gitignore}
written before the first commit would have avoided an afternoon of history
surgery and 2.4 GB of accumulated virtual environments.

% ===================== 10. CONCLUSION ======================================
\chapter{Conclusion}
\label{ch:conclusion}

\section{Summary}

This project delivered a working screenshot triage system for CoMas comprising
a hybrid detection pipeline, two synthetic data generators, a reviewer
application and an evaluation toolchain, supported by 42 automated tests.

Assessed against the objectives of Section~\ref{sec:objectives}:

\textbf{Objective 1 (dataset)} --- partially met. A labelled dataset was
assembled, but synthetically rather than from real CoMas screenshots. The
proposal's stated fallback became the primary strategy.

\textbf{Objective 2 (classifier)} --- met. A ResNet50 classifier was designed,
trained with transfer learning, and integrated into the application.

\textbf{Objective 3 (evaluation)} --- met on synthetic data, not met on real
data. Precision, recall, F1 and AUROC are reported, along with a per-variant
breakdown and a resolution ablation. No quantitative evaluation on real
screenshots was possible.

\textbf{Objective 4 (inference pipeline)} --- met and exceeded. The delivered
application provides ranked review, adjustable sensitivity, triage state and
evidence export.

The most substantive technical finding is that the two visual manifestations of
an AI coding assistant require different detection methods, and that the harder
of the two is resolution-limited in a way that can be measured and quantified
rather than merely asserted.

\section{Limitations}
\label{sec:limitations}

\textbf{Synthetic evaluation only.} The principal limitation, discussed in
Section~\ref{sec:validity}.

\textbf{Ghost-text ceiling.} At 0.856 ghost-vs-clean AUROC the model ranks
ghost text well but does not separate it cleanly. The shipped operating point
recovers 95\% of ghost-text screenshots at the cost of surfacing most clean
editors for review.

\textbf{Keyword maintenance.} Interface strings change with vendor updates. The
keyword list should be reviewed each term, and a strong argument for retaining
the CNN is that it does not share this failure mode.

\textbf{Two tools only.} Only Copilot and Cursor are covered. Other assistants
would be missed by the keyword detector, though the CNN may generalise to
visually similar panels.

\textbf{Single-screenshot analysis.} Each screenshot is judged independently.
Session-level patterns --- a student whose captures show an assistant
repeatedly across an examination --- are not modelled.

\section{Future Work}
\label{sec:future}

\textbf{Obtain real data.} Everything else is secondary. A labelled corpus of
real CoMas screenshots would allow genuine evaluation and would very likely
improve the model, since real screenshots contain variation no generator
reproduces.

\textbf{Tiled inference for ghost text.} Rather than downscaling a whole
screenshot to 768 pixels, divide it into overlapping tiles at native
resolution, score each, and take the maximum. Ghost text would then be
evaluated at full resolution rather than compressed by a factor of four. The
resolution ablation gives direct evidence that this is the most promising
avenue.

\textbf{Session-level analysis.} Filenames already encode session identifiers.
Aggregating scores per student and presenting a timeline would match how
proctors actually think --- about students, not individual images --- and would
allow weak per-image evidence to accumulate into a strong per-session signal.

\textbf{Localisation.} Highlighting the region responsible for a detection,
via attention maps or the tiling above, would let a reviewer verify a flag at a
glance rather than searching the screenshot for what triggered it.

\textbf{Persisting triage state.} Triage decisions are currently held in memory
and lost when the application closes. Persisting them would support review
sessions spanning multiple sittings.

% ===================== REFERENCES ==========================================
\begin{thebibliography}{9}
\addcontentsline{toc}{chapter}{References}

\bibitem{resnet}
K. He, X. Zhang, S. Ren, and J. Sun,
``Deep Residual Learning for Image Recognition,''
\emph{Proceedings of the IEEE Conference on Computer Vision and Pattern
Recognition (CVPR)}, 2016, pp. 770--778.

\bibitem{imagenet}
J. Deng, W. Dong, R. Socher, L.-J. Li, K. Li, and L. Fei-Fei,
``ImageNet: A Large-Scale Hierarchical Image Database,''
\emph{IEEE Conference on Computer Vision and Pattern Recognition}, 2009.

\bibitem{pytorch}
A. Paszke et al.,
``PyTorch: An Imperative Style, High-Performance Deep Learning Library,''
\emph{Advances in Neural Information Processing Systems 32}, 2019.

\bibitem{tesseract}
R. Smith,
``An Overview of the Tesseract OCR Engine,''
\emph{Ninth International Conference on Document Analysis and Recognition
(ICDAR)}, 2007, pp. 629--633.

\bibitem{adamw}
I. Loshchilov and F. Hutter,
``Decoupled Weight Decay Regularization,''
\emph{International Conference on Learning Representations (ICLR)}, 2019.

\bibitem{auprc}
T. Saito and M. Rehmsmeier,
``The Precision-Recall Plot Is More Informative than the ROC Plot When
Evaluating Binary Classifiers on Imbalanced Datasets,''
\emph{PLoS ONE}, vol. 10, no. 3, 2015.

\bibitem{transfer}
J. Yosinski, J. Clune, Y. Bengio, and H. Lipson,
``How Transferable Are Features in Deep Neural Networks?''
\emph{Advances in Neural Information Processing Systems 27}, 2014.

\end{thebibliography}

% ===================== APPENDICES ==========================================
\appendix

\chapter{Repository Structure}

\begin{lstlisting}[language=,caption={Project layout. Generated data and
checkpoints are excluded from version control.}]
Honours Project/
  comas/                        main package
    __init__.py
    brightspace.py              Brightspace detector
    config.py                   typed YAML configuration
    copilot.py                  VS Code detector, CNN scorer
    data.py                     dataset, transforms, splitting
    evaluate.py                 metric computation
    gui.py                      desktop reviewer interface
    infer.py                    command-line inference
    model.py                    ResNet50 and fusion models
    ocr.py                      Tesseract wrapper and cache
    realpairs.py                real screenshot utilities
    synthetic.py                VS Code screenshot generator
    synthetic_brightspace.py    browser screenshot generator
    train.py                    training loop
    variant_report.py           per-variant evaluation
    assets/carleton_logo.png
  tests/                        42 automated tests
  data_vscode/                  generated training data
  data_vscode_holdout/          generated evaluation data
  data_brightspace/             generated browser data
  checkpoints/                  trained model weights
  cache/ocr.json                OCR cache
  config.yaml                   configuration
  requirements.txt
  README.md
\end{lstlisting}

\chapter{Configuration Reference}

\begin{lstlisting}[language=,caption={\code{config.yaml} as used for the final
768-pixel model.}]
data:
  root: data_vscode/
  img_size: 768
  num_workers: 4
  batch_size: 8
  val_frac: 0.15
  test_frac: 0.15
  seed: 42

model:
  use_ocr: false
  cnn_backbone: resnet50
  text_encoder: sentence-transformers/all-MiniLM-L6-v2
  fusion_hidden: 256
  dropout: 0.3

train:
  epochs: 10
  lr: 1.0e-4
  weight_decay: 1.0e-4
  pos_weight_auto: true
  pos_weight_override: null

paths:
  checkpoint: checkpoints/vscode_768.pt
  ocr_cache: cache/ocr.json

logging:
  level: INFO
\end{lstlisting}

\chapter{Instructions to Run and Test}

\section{Environment}

Python 3.11 is required. On macOS the system Python bundled with Xcode must
\emph{not} be used, as its Tk 8.5 has rendering defects that prevent the
interface from displaying correctly (Section~\ref{sec:tkbug}).

\begin{lstlisting}[language=bash]
brew install python@3.11 python-tk@3.11 tesseract
python3.11 -m venv .venv
source .venv/bin/activate
pip install -r requirements.txt
python3 -c "import tkinter; print(tkinter.TkVersion)"   # must print 8.6
\end{lstlisting}

\section{Running the Application}

\begin{lstlisting}[language=bash]
python3 -m comas.gui
\end{lstlisting}

Select a detection mode, upload screenshots, and review the ranked results.
Keyboard shortcuts in the review screen: arrow keys navigate, \code{r} marks
reviewed, \code{d} dismisses, \code{c} clears, \code{e} exports.

\section{Regenerating Data and Retraining}

\begin{lstlisting}[language=bash]
# generate datasets
python3 -m comas.synthetic --out data_vscode --n-active 200 --n-clean 200
python3 -m comas.synthetic_brightspace --out data_brightspace
cp data_brightspace/on_quiz/*.png data_brightspace/left_quiz/*.png \
   data_vscode/no_copilot/

# train
python3 -m comas.train

# held-out evaluation
python3 -m comas.synthetic --out data_vscode_holdout \
    --n-active 100 --n-clean 100 --sessions 10 --seed 777
python3 -m comas.variant_report --data data_vscode_holdout \
    --split all --threshold 0.20
\end{lstlisting}

\section{Running the Tests}

\begin{lstlisting}[language=bash]
python3 -m pytest tests/ -q
\end{lstlisting}

Expected result: 42 passed. Tests requiring Tkinter are skipped automatically
in environments without a display.

\section{Command-Line Detection}

\begin{lstlisting}[language=bash]
python3 -m comas.copilot --folder path/to/screenshots --out ranked.csv
python3 -m comas.brightspace --folder path/to/screenshots --out bs.csv
\end{lstlisting}

\end{document}
